\section{The Mystery of Birth}

\begin{quotex}
To allow such traditional views to act on oneself, until they awaken precise interior forces, means to progress on the way that leads to a spiritual level from which the significance of life constitutes something absolutely different from what it means for the rest of men: a significance of clarity, absolute force, and incomparable certitude.
\flright{\textsc{Julius Evola}}
\end{quotex}


In \emph{Sintesi di dottrina della razza}, Evola addresses the mystery of birth and how it relates to one's destiny. After discussing the notions of the races of the soul and spirit, he raises the question and addresses some objections to his view:

\begin{quotex}
What blame does a man have if he was born in one race and not in another? Is he perhaps responsible for the fact that his parents and ancestors are Aryans, Jews, blacks or Amerindian? Has he perhaps willed all that? With your theory of race you remain, in spite of it all, in a purely naturalistic point of view. You turn a natural event into a destiny and construct above your system, instead of bringing attention above all to those values in which it can truly enter into play, and be considered imputable, as a human responsibility.
\end{quotex}


He regards these objections not as weird or artificial, but of a real importance to avoid reverting to materialistic and collectivist degradations of the doctrine. The traditional point of view stresses the values of the person over any of its manifestations.

Evola asserts that only pre-christian conceptions of man and life can satisfy his requirements. We prefer to deal with concepts rather than labels, so we need to reformulate Evola's argument somewhat. In a subsequent post, we will show how his views challenge both the traditional right and its modern alternatives. Evola describes the Christian solution to the mystery of birth this way:

\begin{quotex}
The Church had to reject the idea of pre-existence that the preceding traditions always recognized: that is, it denied that the spiritual nucleus of the personality pre-exists terrestrial birth and, naturally, conception. In Christian theology things in this respect are not so simple, as far as this denial could be believed. Nevertheless, the fundamental Christian view is that each human soul is unique and is created by God from nothing at the moment in which it is breathed into a body or embryo read to receive it. That a man is born in one race rather than another becomes then a theological mystery: ``God willed it'' and, often, it is admitted that the divine will is inscrutable.
\end{quotex}


In that paragraph, there is something true, something false, and something ambiguous. It is true that the human being exists from the moment of conception, so, here, Evola is more Catholic than the pope. It is false that Catholic doctrine teaches that the soul is infused into the embryo at some point subsequent to conception. That is the Cartesian view that presumes a duality of soul and body. What is ambiguous is the meaning of the pre-existence of the soul. The common misunderstanding is that the soul exists in time prior to conception. This is the view of Plato of all systems deriving from Theosophy, and even popular views of oriental religions such as Hinduism or Buddhism.

To clear up the latter point first, Evola makes it clear that this pre-existence cannot be conceived in a temporal sense, since the human Self is transcendent to the world of space and time. The Thomist view, then, is not as dissimilar as Evola would like. The soul is the form of the body, which cannot exist without the soul. So the soul is pre-existent in the ontological sense and the physical aspects of the body are the reflection of soul qualities. This is really Evola's point, and should be the Catholic viewpoint, had the understanding brought by Thomas Aquinas allowed to be further developed instead of being petrified at a certain point. Evola explains his doctrine:

\begin{quotationx}
The view of ancient Aryan humanity was totally different and it alone permits us to overcome the objections indicated... According to that view, birth is not by chance nor a fact willed by God; neither fidelity in regards to one's own nature signifies a passivity, but rather testifies to the more or less clear knowledge of a profound connection of one's Self with something transcendental and otherworldly, to be able to act in a transfiguring way. This is the essence of the doctrine of karma and dharma, a doctrine that must not be confused with the idea of reincarnation.

It is necessary to distinguish a double heredity. That which stands before the individual in the temporal non-transcendental, sense is the heredity of parents, the gens, the race, a certain civilization, caste, etc., therefore, more or less, everything that commonly is meant when speaking of heredity. But all that does not exhaust the spiritual reality of the individual, as materialism and historicism claim: as determining and essential it must rather be considered an intervention from above, a principle assuming and utilizing, as its material of expression and incarnation, everything that this heredity has gathered with its laws and determinism. Moreover, it must be thought that the historical-biological heredity of a given line is chosen and assumed when it can be approximately worth as a type of expressions analogous to a transcendental heredity.

That is why in every being two heredities meet and join each other, the one terrestrial and historic, in good measure positively detectable i.e., known to science (tr) , the other spiritual, otherworldly. To establish the ties between the two, is therefore to determine the synthesis that defines a given human nature, and intervenes in an event represented in various traditions with various symbols... Basically, as noted, a type of law of elective affinity acts here. Wishing to clarify with some applications, we will say that a man is not of a given gender, race, or caste because he is born that by chance or through the will of God, or through a mechanism of natural causes, but vice versa. He was born that way because he already was a man or women, of one race or another, etc., naturally in an analogical sense in the sense of an inclination of vocation or deliberation transcending use, through lack of adequate concepts, we can present only through its effects. In a certain way, one has therefore the interference of horizontal lines and vertical lines of an earthly heredity and non-earthly. At the point of crossing, according to traditional teachings, birth occurs, or better said, the conception of a new being, an incarnation.

Race, caste, etc., therefore exist in the spirit before manifesting in terrestrial and historical existence. Diversity has its origin from above, and what refers to it on earth is only a reflection and symbol. As one willed to be on the basis of a primordial nature or a transcendental decision, so one is. It is not birth that determines nature, but vice versa; it is nature, in its widest sense, because common words are inadequate, that determines birth.
\end{quotationx}


It is this point of crossing that accidental, or contingent, qualities intersect with transcendent, or essential, qualities. Clearly, the modern conception ignores the transcendent, by its nominalism and denial of formal causes, hence every quality is accidental. The Aristotelian-Thomist scheme agrees in principle, yet is limited by its commitment to the genus and species distinction, so it is forced to define various types of accidents. We will address these issues subsequently, after we translate another section on the impact that Evola's conception has on ethical theory. But first we will repeat some quotes from Plotinus and Plato that Evola relies upon.

\begin{quotex}
The general plane is unique, but divides itself in diverse parts, in a way that in everything there are distinct places, some more pleasant, others less so. And souls, also unequal, come to inhabit these distinct places that correspond to their own diversity. In such a way, everything is arranged and the differences in their situations correspond to the inequality of the souls.
\flright{\textsc{Plotinus}}
\end{quotex}


\begin{quotex}
It is not the demon that chooses you, but you yourself chose your demon. You yourself choose the destiny of that life to which you are irremediably connected.
\flright{\textsc{Plato}}
\end{quotex}

\paragraph{Mystery of Birth}


For the presentation of the directive principles of this part of race theory it is however necessary to make a few remarks related to the problem of birth, and a final clarification of what has been said about heredity.

Even when we have worked out all the principal objections against the doctrine of race that are made from an immediate, practical or intellectual point of view, whether in good or bad faith, there remains one that is insuperable but decisive. They can say: good, everything affirmed is correct. But with everything taken into account, what fault does a man have if he was born into one race and not in another? Is he perhaps responsible for the fact that his parents and ancestors are Aryans, Jews, blacks or Amerindian? Has he perhaps willed all that? With your theory of race you remain, in spite of everything, in a purely naturalistic point of view. You turn a natural event into a destiny and you build above your system, instead of paying attention above all to those values in which human responsibility can truly come into play, and be considered imputable.

This is, in some way, the ultima ratio of the adversaries of race theory. And it must be conceded that such an objection is not artificial and farfetched, but of a real importance, if it does not support the materialist and collectivist deterioration of the doctrine in word and is posed instead from the traditional point of view, which always emphasizes the values of the personality. Nevertheless, to consider that objection certainly means confronting the problem of birth. From a higher, spiritual point of view, the justification of the race idea depends on the problem of birth and its solutions.

In this regard, arriving at some firm points is still rather difficult, as long as it remains in the ambit of views introduced by the coming of Christianity into the West. And that is not by chance: race and super-race, cult of blood, arianity, etc., are all concepts that were formed and affirmed essentially in pre-Christian civilizations. It is therefore necessary to look in such traditions and in their knowledge for the elements of the solution of the problems, that the recovery of those ideas today are going to raise. Every reference in more recent conceptions of man and life can only furnish us incomplete and often inadequate points of view.

Thus we must not be surprised that the problem of birth is remarkably obscure in the Christian vision of the world. For precise and certainly not arbitrary reasons, that we cannot discuss here, the Church had to reject the idea of pre-existence that earlier traditions had always recognized: that is, it has denied that the spiritual center of the personality pre-exists terrestrial birth and, certainly, even conception. In Christian theology, things in this regard are not as simple, as this denial can be believed. Nevertheless, the fundamental Christian view is that each human soul is unique and is created by God from nothing at the moment in which it is breathed into a body or human embryo ready to receive it. That a man is born into one race rather than another therefore becomes a theological mystery: ``God willed it'' and it is usually admitted that the divine will is inscrutable.

The view of ancient Aryan humanity was totally different and it alone can overcome the indicated objections. For a complete explanation, we must again remind the reader of our work Revolt against the Modern World. Recapping, here we limit ourselves to say that according to that view, birth is never by chance, nor a fact willed by God; nor does fidelity in respect to one's own nature signify passivity, but rather it gives evidence to the more or less clear consciousness of a deep connection of one's own Self with something transcendent and otherworldly, to the extent of acting in a transfiguring way. This is the essence of the doctrines of karma and dharma, which must not be confused with the idea of reincarnation. As we have shown in other works, the theory of reincarnation is either a conception foreign to ``Aryan spirituality, essentially characteristic of pre-aryan, telluric, matriarchal cycles of civilization, or it is the effect of misunderstandings and deformations that certain traditional views have undergone in some contemporary theosophical circles. And if, in the traditional world, even in the Aryan, seemingly clear evidence in favor of the belief in reincarnation, in reality, it is only about symbolic forms that a superior knowledge had to disguise in regard to the people and the uninitiated.

In any case, for the problem that concerns us here, it is necessary to refer, not to reincarnation, but to the doctrine that the human Self, as I having my own given nature, would be the effect, the production, the mode of appearance under certain conditions of existence, of a spiritual entity that preexists and transcends it. And since everything that is time, sooner or later, is just something inherent in the human condition, so strictly speaking we cannot speak of anything preexisting or antecedent in the temporal sense.

We are getting into a rather difficult field, precisely because conceptions and expressions cannot be applied to it that we are formed in the existence of this world and that, applied to a different reality, they can easily lead to distortions and deformations. We say, in any case, that It is necessary to discern a double heredity. What stands before the individual, in the temporal, non-transcendental sense, is the heredity of parents, the folk, the race, a certain civilization, caste, etc., therefore, more or less, everything that commonly is meant when speaking of heredity. But even that does not exhaust the spiritual reality of the individual, as materialism and historicism claim: rather we must consider an intervention from above as determining and essential, a principle assuming and utilizing everything that this heredity has accumulated as its material of expression and incarnation, with its laws and its determinisms. Moreover, it must be thought that the historical and biological heredity of a given line is chosen and assumed when it can be approximately be considered as a kind of expressions analogous to transcendental heredity.

That is why two heredities meet and join each other in every being, the one terrestrial and historical, positively detectable to a great degree i.e., known to science (tr) , the other spiritual and otherworldly. To establish the connection between the two, and thereby determine the synthesis that defines a given human nature, an event intervenes, represented in various traditions with various symbols, which we cannot examine here. Basically, as noted, a type of law of elective affinity operates. In order to clarify with some applications, we can say that a man is not of a given gender, race, or caste because he is born that way by chance or through the will of God, or through a mechanism of natural causes, but vice versa. He was born that way because he already was a man or women, of one race or another, etc., naturally in an analogical sense in the sense of an inclination of a transcendent vocation or deliberation, which we, because of the lack of adequate concepts, can indicate only through its effects. In a certain way, we therefore have the interference of horizontal lines and vertical lines of an earthly heredity and an otherworldly. At the point of crossing, according to traditional teachings, birth occurs, or better said, the conception of a new being, an incarnation.

Race, caste, etc., therefore exist in spirit before manifesting in terrestrial and historical existence. Diversity has its origin from above, and what relates to it on earth is only a reflection and symbol. As you have willed to be, based on a primordial nature or a transcendental decision, so you are. It is not birth that determines nature, but vice versa; it is nature, in its widest sense, because common words are inadequate, that determines birth.\footnote{For further reading, see ``The Doctrine of the Castes'' in \emph{Revolt against the modern World}, ``Man and Woman'' in \emph{Revolt against the modern World}, \textit{Life of the Spirit} (\url{https://gornahoor.net/?p=2163}), \textit{Law of Attraction} (\url{https://gornahoor.net/?p=1664}).}

\paragraph{The Mystery of Birth Explained}

\begin{quotex}
The general plane is unique, but divides itself in diverse parts, in a way that in everything there are distinct places, some more pleasant, others less so. And souls, also unequal, come to inhabit these distinct places that correspond to their own diversity. In such a way, everything is arranged and the differences in their situations correspond to the inequality of the souls.
\flright{\textsc{Plotinus}}

It is not the demon that chooses you, but you yourself chose your demon. You yourself choose the destiny of that life to which you are irremediably connected.
\flright{\textsc{Plato}}
\end{quotex}

After the quotes from Plotinus and Plato, \textbf{Julius Evola} next discusses various supernatural beings from Tradition, which Evola accepts as real entities and essential for his ideas on spiritual races. However, there is no reason to restrict their influence to antiquity; they are present and fully active today.

\begin{quotex}
This expression ``demon'', from Plato is particularly interesting for us, since here the concept of demon has nothing to do with the Christian idea of a wicked entity, but instead has a close relation to the deepest forces of races, both of the soul and of the body. Even here, we cannot delve into the traditional doctrine in regards, but only to recall that, likewise, the demon, the lares, the penates, the ``doppelganger'' \emph{doppio}, which is a synonym for ``subtle body'', are entities that in antiquity interfered in and reflected the precise knowledge of the true roots of the differentiation of bloodlines, \emph{gentes}, and ultimately the individuals themselves on the basis of a totalitarian vision of the world, including the invisible and the visible, and not of that mutilation of the moderns who only know material and psychological processes. From such testimonies which could be multiplied with reference to the traditions of all peoples, the idea of transcendental, or vertical, heredity therefore is confirmed, and of the choice that on the base of analogous correspondences, determines the connection of it to a horizontal, historical, biological heredity. The consequences of all the in regards to the justification of the racial ideas are clearly visible.
\end{quotex}


Since we have already discussed the notions of Angels and Demons\footnote{\url{https://gornahoor.net/?p=2995}}, we will move onto the other supernatural beings. The lares\footnote{\url{https://en.wikipedia.org/wiki/Lares}} are protecting spirits of dead ancestors. They kept watch over the affairs of the family. Many pious people today are aware of the watchful presence of their ancestors. In addition, there are lares to protect broader domains such as neighborhoods, cities, states, livestock, and so on. Obviously, we still know these spirits, but they are now called patron saints or guardian angels. The principalities are the guardian angels for the various nations and countries, as accepted by Muslims and Traditional Christians, as well as the ancient pagans.

The penates\footnote{\url{https://en.wikipedia.org/wiki/Di_Penates}} more or less serve the same function as the lares. Since Evola accepts demons, lares, and penates as real beings, we should expect their presence to be made known, at least to those of the proper disposition and sensitivity. That they are called today the spirits of one's ancestors, patron saints, or angels makes no difference for their reality.

By ``double''\footnote{\url{https://en.wikipedia.org/wiki/Doppelg\%C3\%A4nger}}, doppio, or doppelganger, Evola means the subtle body, i.e., the soul. In ancient Egypt this was called Ka. In the \emph{Ancient City}, \textbf{Fustel de Coulanges} describes the state of the soul after death, as related in Hindu, Greek, and Roman traditions. Since we have previously discussed this on several occasions, and in the future we want to compare it to the state of the souls in Limbo according to Dante, we will move on.

So it should be absolutely clear now what Evola means by the ``third dimension'' of history. The modern mind can only see material or historical processes, which correspond to gross and subtle manifestation. However, the vertical dimension includes a whole array of beings who are involved in the creation of nations, peoples, and individuals. The actions of these beings are invisible to the modern mind which can only attribute them to material and biological causes.

Given this, we are in a position to better understand the existence of man in the light of Tradition. As Evola points out, the difference between Essence and Existence is a matter of degree. Here we understand man's Essence as who he is in himself, his Spirit, his karma (what he carries into the world) and his dharma (his task or destiny in the world). His Existence, then, is the meeting of this Essence interacting with the horizontal forces at the proper time and space suitable to him, by the Law of Elective Affinity. This is not a bipolar process between Essence and Existence, but rather an emanation involving higher beings and influences. If you look at the diagram from \textit{Spiritual Beings}\footnote{\url{https://gornahoor.net/?p=1941}}, this process takes the being through the angelic influences which determine his race, nation, family, and so on, then through the astral and planetary influences, and finally through the four elements. This is the birth of man in the vertical direction.

\flrightit{Posted on 2012-04-05 by Aeneas}

\begin{center}* * *\end{center}

\begin{footnotesize}
\begin{sffamily}
\texttt{Temenus on 2012-04-06 at 06:39 said:}

The Plato quote is taken from RepubliC,Book 10, 617e. Here is an other translation of the whole passage that I found (I think that the translation of greek ``daimon'' as divinity and deity can be misleading) : 617e No divinity shall cast lots for you, but you shall choose your own deity. Let him to whom falls the first lot first select a life to which he shall cleave of necessity. But virtue has no master over her,and each shall have more or less of her as he honors her or does her despite. The blame is his who chooses: God is blameless. ``So saying, the prophet flung the lots out among them all, and each took up the lot that fell by his side, except himself; him they did not permit.And whoever took up a lot saw plainly what number he had drawn.''

\url{http://www.perseus.tufts.edu/hopper/text?doc=plat.+rep.+10.617e}


\hfill

\texttt{golgonooza on 2012-04-09 at 15:39 said:}

In `Lazarus, Come Forth' Valentin Tomberg discusses the second miracle of John's Gospel, the healing of the nobleman's son, and the sixth day of creation in terms of two types of heredity -- `vertical' and `horizontal'. The Fall is seen by him as the substituting of the vertical heredity of creation in image and likeness of God for the horizontal heredity of ancestral traits, original sin etc. The healing of the nobleman's son placed him into direct relationship with Christ allowing the heredity of the father to be subsumed into the `vertical' relationship. This seemed to me to be very similar to what Evola is saying here -- am I right?

\hfill

\texttt{Cologero on 2012-04-09 at 17:46 said:}

Thanks for reminding me of that passage. Yes, it is the same understanding although Evola would avoid using the words ``God'' or ``Christ''. Of course, Tomberg is drawing on the same tradition as Evola from Hermes to Plotinus. However, Evola is too polemical and eager to see differences, where Tomberg and Guido de Giorgio see continuity and a deeper development.

\hfill

\texttt{Logres on 2012-04-09 at 22:56 said:}

Ortega y Gasset cites the Chinese custom of worthy deeds ennobling the ancestor. Transferred to a spiritual plane (vertical plane), is this also the same thought?

\hfill

\texttt{Angolmois on 2012-04-28 at 13:55 said:}

Ancestors and God should be brought together in one's worship, not to separate them, but of course the Divine, the vertical should have a pre-eminence. (For the record, I'm not implying that this would be recommended here, just mentioning it.)

What is the most presumptuous and hypocritical comment one hears alot in favouring multio-racialism and multi-culturalism nowadays is that Aryans and Non-Aryans come in all colours and nationalities. This has never been in doubt anywhere in world's traditions, but somehow these modern idiot's seem to think they're wiser and that all the world Aryans should mix with each other everywhere and then we'd have the grey-black-brown globalist utopia they hunger for. Spit and shame on them.

\hfill

\texttt{Cologero on 2012-04-28 at 17:37 said:}

Angolmois, I think you mean that as the One is manifested, it becomes more and more differentiated, not less. The modern mind, denying the supernatural \emph{de facto}, even if they confess the opposite, believes instead that the One must appear on the material plane. Hence, they deny all differentiation and hierarchy in the face of all evidence.

40,000 Norwegians singing together a communist inspired song for peace changes nothing. Moltke said: ``universal peace is only a dream, and not even a beautiful dream.''

There are still some who are willing to combat; they may even have an inkling of their enemy (the negative pole). But they will not be effective unless they know precisely what they are fighting for (the positive pole).

\hfill

\texttt{Angolmois on 2012-04-29 at 03:08 said:}

Yes, that is what I meant with my comment. The differentiation of the divine and its multiple manifestations in the form of different races, colours etc. should and could be a compliment for representatives of different races, nationalities and traditions -- separate yet unified -- but instead the modernists have made it an insult with their dis-honesty and linguistic / semantic plays. For example, ``a negro'' is a four letter curse word nowadays, whereas it was a compliment and a honorary title for the representatives of their race before. It has nothing to do with ``hate speech'' or anything like that for sure.

\hfill

\texttt{Timotheus Lutz on 2013-01-08 at 22:52 said:}

Evola wrote: `In a certain way, one has therefore the interference of horizontal lines and vertical lines of an earthly heredity and non-earthly.'

I wonder if Francis Parker Yockey, noted American Axis-sympathizer, read and was influenced by the Synthesis. The characterization of spiritual and physical race as horizontal and vertical appeared in his book Imperium, although if I'm remembering right, the words were switched.

\end{sffamily}
\end{footnotesize}

